\chapter{Podsumowanie i wnioski}

Kwantowe maszyny cieplne to aktywna dziedzina badań naukowych. Mają one zastosowanie zarówno jako silniki jak i chłodziarki oraz pompy ciepła.

Cykl Otto wykorzystujący oscylator harmoniczny jako czynnik roboczy w niskich temperaturach wykazuje w punkcie maksymalnej mocy sprawność wyższą niż w przypadku klasycznych cykli Carnota, Otto oraz Stirlinga dla dowolnego zakresu stosunku temperatur zbiorników cieplnych, z którymi jest w kontakcie. Ten sam cykl wykorzystujący układ o spinie \(\inv{2}\) ma sprawność wyższą od silników klasycznych jedynie w pewnym zakresie stosunku temperatur. W wyższych temperaturach sprawność pierwszego z nich pokrywa się z silnikami klasycznymi, ale dla drugiego jest ona zawsze niższa. 

Cykl Stirlinga z regeneratorem w warunkach równowagowych pracujący na układzie o spinie \(\inv{2}\) może teoretycznie osiagnąć sprawności wyższe niż ta idealnego silnika Carnota. Jest to wynik wyidealizowany i taki regenerator nie może pracować bez zewnętrznego źródła energii. Ten sam cykl bez regeneratora ograniczony jest limitem idealnego silnika Carnota. Dla cyklu wykorzystującego oscylator harmoniczny sprawność jest zawsze niższa od tej silnika Carnota, ale regenerator zwiększa sprawność. Wydajne działanie regeneratora jest ważne w celu osiągnięcia wysokiej sprawności cyklu.